\section{Method}
\label{sec:method}

\method{} converts a user's question into an interactive visual image through
a pipeline that decouples \emph{content structure} (decided deterministically
by the system) from \emph{visual placement} (discovered after image
generation). This section describes the end-to-end pipeline
(Section~\ref{sec:method_pipeline}), the structured answer normalization and
layout planning (Section~\ref{sec:method_structure}), the two-pass vision
alignment that grounds hotspots to real image content
(Section~\ref{sec:method_alignment}), and the interactive hotspot layer and
per-region follow-up threads (Section~\ref{sec:method_interaction}).

\begin{figure}[t!]
    \centering
    \includegraphics[width=\linewidth]{model.pdf}
    \caption{\textbf{\method{} pipeline.} A user question yields a raw LLM
    answer, which is normalized into a structured visual specification and a
    normalized layout. An image prompt is synthesized and the image is
    generated. In the second pass, a visual grounding model localizes where
    each structured module actually rendered; hotspots are placed at the
    grounded coordinates and become clickable regions with detail panels and
    follow-up threads.}
    \label{fig:architecture}
    \vspace{-0.6em}
\end{figure}

\subsection{End-to-End Pipeline}
\label{sec:method_pipeline}

Given a user question $q$, \method{} produces an interactive visual image
$\mathcal{I}$ with a set of clickable hotspots $\mathcal{H}$. The pipeline
proceeds in two passes (Figure~\ref{fig:architecture}):

\paragraph{Pass 1: Structure, layout, and image generation.}
\begin{enumerate}[leftmargin=*,topsep=2pt,itemsep=1pt]
    \item \textbf{Raw answer.} Obtain a long-form LLM answer $a$ to $q$ via a
    text model (or a deterministic mock for development).
    \item \textbf{Normalization.} Normalize $a$ into a structured visual
    specification $\mathcal{S}$---a set of modules $\{m_i\}$ and auxiliary
    modules, each with a title, detail text, region prompt, and metadata
    (Section~\ref{sec:method_structure}).
    \item \textbf{Layout planning.} Plan a layout $\mathcal{L}$ that assigns
    each module a region $r_i$ with normalized bounds in $[0,1]^4$.
    \item \textbf{Image prompt synthesis.} Derive an image prompt from
    $\mathcal{S}$ and $\mathcal{L}$, encoding the visual style and the
    intended region labels and positions.
    \item \textbf{Image generation.} Generate the image via a text-to-image
    model, obtaining $\mathcal{I}$ and its real pixel dimensions
    $(W, H)$ parsed from the image header (not the request, since generators
    do not guarantee exact sizes).
\end{enumerate}

\paragraph{Pass 2: Vision alignment and hotspot derivation.}
\begin{enumerate}[leftmargin=*,topsep=2pt,itemsep=1pt,resume]
    \item \textbf{Visual grounding.} Send $\mathcal{I}$, $(W,H)$, and the
    modules with their planned bounds to a visual grounding model, which
    returns a grounded box $b_i$ for each module where it actually rendered
    (Section~\ref{sec:method_alignment}).
    \item \textbf{Mask refinement (optional).} Optionally refine each $b_i$
    with a segmentation model (SAM3) to obtain a pixel mask for tighter
    previews.
    \item \textbf{Hotspot derivation.} Copy the grounded bounds into hotspot
    geometry $h_i = (x_i, y_i, w_i, h_i)$ and render transparent, clickable
    hotspots over $\mathcal{I}$ (Section~\ref{sec:method_interaction}).
\end{enumerate}

The key design principle is that \method{} \textbf{does not assume the image
generator follows the planned layout}. Instead of requiring precise spatial
control---which diffusion models do not reliably provide~\citep{feng2023layoutgpt}
---the system \emph{generates first, then grounds}. This makes hotspot
alignment robust to the layout imprecision of generative models.

\subsection{Structured Answer Normalization and Layout}
\label{sec:method_structure}

\paragraph{Visual specification.}
The raw answer $a$ is parsed (via a structured-output LLM call) into a visual
specification $\mathcal{S}$ with a visual mode (infographic, map, scene, or
poster), a summary, and a list of modules. Each module $m_i$ carries:
\textit{title} (the hotspot label), \textit{detail} (the explanation shown in
the detail panel), \textit{imageText} (short on-image caption),
\textit{regionPrompt} (a locator phrase for visual grounding), and metadata
(\textit{regionKind}, \textit{iconHint}, \textit{maskPolicy}, locator queries).
Because LLMs sometimes echo instruction-style text into the user-facing
\textit{detail} field, we sanitize it: a cleaning function strips
visual-system vocabulary and meta-instruction phrases clause-by-clause,
falling back to \textit{imageText}/\textit{title} when nothing usable survives.
When the LLM returns no valid structure, \method{} falls back to a library of
content-grounded mock specifications (Section~\ref{sec:experiments} lists the
topics), ensuring the system always produces a valid visual answer.

\paragraph{Layout planning.}
A layout planner assigns each module a region $r_i$ with normalized bounds
$\in [0,1]^4$ relative to the image canvas. Layouts are organized by
family---grid, flow (swimlane), compare, hub, timeline, matrix, and
freeform---selected by the visual mode and relation type. For maps, a
semantic-slot catalog assigns regions by keyword matching (e.g., ``entrance''
$\to$ bottom-center, ``core area'' $\to$ center); for infographics, index-based
templates place modules in flow or grid order. These \emph{planned} bounds
serve as priors for the grounding step and as a fallback when grounding fails,
but they are \emph{not} assumed to match the generated image exactly.

\subsection{Two-Pass Vision Alignment}
\label{sec:method_alignment}

The core technical contribution is aligning the structured regions to the
generated image. \method{} sends the image, its real dimensions, the modules,
and their planned bounds to a vision backend that runs a grounding pipeline:

\paragraph{Grounding provider.}
The grounding step is provider-agnostic: \method{} ships a JSONL worker
interface and any model that returns a normalized bounding box per locator
phrase can be plugged in. Our reference implementation supports two
interchangeable backends. The default is a dedicated referring-localization
model (\textbf{LocateAnything-3B}~\citep{nvidia2025locateanything}), launched
through the \texttt{locateanything\_worker.py} worker. As an alternative we
also support a general vision-language model used in a structured-output
grounding mode (\textbf{MiMo-Vision}~\citep{xiaomi2025mimovl}); this is
the path actually used by several of our released demos
(Section~\ref{sec:experiments_benchmark}). For each module the backend is
given the image and one or more locator phrases built from the module's
title, region prompt, and semantic hints. It outputs a bounding box per
module in normalized coordinates, scored against the planned bounds to
select the best candidate (with optional crop-and-re-ground for low-confidence
cases). This is a \emph{referring localization} step: we are localizing
``where the module labeled $X$ actually rendered,'' not detecting a
predefined object class.

\paragraph{Fallbacks.}
When the grounding backend rejects a module or is unavailable, \method{}
falls back along a graded chain: a general VLM-based grounding pass, then
local OCR, and finally the planned bounds. Each module is tagged with its
alignment source (e.g., \textit{locateanything}, \textit{mimo-vision},
\textit{sam3-refined-planned}, \textit{planned-fallback},
\textit{vision-low-confidence}) so the system can report alignment quality
per region rather than as a single aggregate score.

\paragraph{Robustness to partial failure.}
A single missing or low-confidence grounding must not corrupt the whole
result. We reject groundings that look like header strips or cross-panel
artifacts (falling back to planned bounds for those) while keeping
far-from-planned groundings that still localize the correct region (tagged
low-confidence). Modules entirely absent from the grounding response retain
their planned bounds and are flagged, rather than discarding the alignments of
all correctly-grounded modules. Each grounded box is repaired for
clickability---enforcing a minimum area and safe margin---before becoming a
hotspot.

\paragraph{SAM3 mask refinement (optional).}
For tighter visual previews, an interactive box-prompted segmentation
model in the SAM family~\citep{kirillov2023sam,ravi2024sam2}---we use the
SAM3 image checkpoint released by Meta\footnote{\url{https://github.com/facebookresearch/sam3}}
through the \texttt{sam3\_worker.py} worker---takes each grounded box as a
prompt and predicts a pixel mask, from which a tightened box and polygon
are derived. The mask is used only for preview rendering (cutout, organic
feathering, soft fade); \emph{click hit-testing always uses the rectangular
box}, keeping interaction robust even when masks are imperfect.

\subsection{Interactive Hotspot Layer and Follow-up Threads}
\label{sec:method_interaction}

\paragraph{Hotspot rendering.}
Each hotspot is a transparent, absolutely-positioned button over the image at
its grounded coordinates. In normal use, hotspots are fully invisible so the
generated image is the only thing the user sees; a hover state provides a soft
visual hint without altering the image. A calibration mode toggles visible
outlines and per-hotspot labels for inspecting and correcting bounds.

\paragraph{Detail panels and previews.}
Clicking a hotspot opens a centered detail panel showing the module's title,
sanitized detail text, a localized preview of the image region, and a
follow-up input. The preview variant is chosen by a strategy module based on
the hotspot shape and mask availability: cutout (checkerboard background),
organic feathering, soft radial fade, or mask-based polygon crop.

\paragraph{Per-region follow-up threads.}
Each hotspot owns an isolated follow-up thread. A follow-up question is sent
to the text model with a \emph{local} context---the original question, the raw
answer, the clicked hotspot's detail, and that thread's history---so the
answer is scoped to the region without re-explaining the whole image. Threads
are persisted alongside the ChatImage, enabling resumption across sessions.

\subsection{Implementation}
\label{sec:method_implementation}

\method{} is a provider-agnostic system with a browser frontend, a local Node.js
backend, and Python model workers. The frontend is vanilla JavaScript (no
framework); the backend proxies text, image, and vision providers and persists
results in SQLite. Vision alignment runs via persistent JSONL workers
(\texttt{locateanything\_worker.py}, \texttt{sam3\_worker.py}) hosting the
grounding and segmentation models. The system runs in a no-key \textit{mock}
mode (deterministic local providers and SVG output) for development and demos,
or in \textit{api} mode proxying real models. All code, demos, and a live
project page are open-source.\footnote{\url{https://github.com/wencanjiang/ChatImage}}
